\section{Conclusions}\label{sec:conclusions}
% P:conclusion-01
The absorber formulation connects reactive ePC-SAFT thermodynamics to a constrained equilibrium-manifold film and countercurrent material and energy balances.
The local film resistance is obtained by quadrature in reactive state space and a scalar interface balance, while the axial column remains a two-point boundary-value problem.
Consistent species amounts, conserved totals, and enthalpy references connect equilibrium, transport, and the thermal response.

% P:conclusion-02
The comparisons relate thermodynamic model choice, parameter sensitivity, film resistance, and computational effort to the same capture and temperature observables.
The local operating study places those dependencies within a specified range of feeds and property inputs.
% INSERT final quantitative conclusions for thermodynamics and parameter sensitivity, transfer, solution cost, and operating response, in that order.
% INSERT the engineering implication supported jointly by capture and axial-temperature comparisons.

% P:conclusion-03
The thermodynamic study connects the Henry-law, eNRTL, and ePC-SAFT driving-force descriptions to capture and axial temperature.
The common-property comparison identifies the response to the selected thermodynamic description, including its specified species and reaction system.
The full-property comparison identifies the additional density and caloric response under the same transport assumptions.
The parameter sensitivity then distinguishes the influence of reaction-enthalpy directions, water solvation, and CO$_2$ dispersion within the selected ePC-SAFT description.
% CHECKLIST-BEGIN: conclusion-1
% P:conclusion-04
\textit{[Insert the principal thermodynamic finding, its quantitative effect, and the parameter-sensitivity finding that qualifies it.]}
% CHECKLIST-END: conclusion-1

% P:conclusion-05
The transfer study compares the enhancement-factor and equilibrium-manifold resistance with the same thermodynamic model and feed conditions.
The interfacial state, flux distribution, and cumulative uptake explain the resulting column response.
% CHECKLIST-BEGIN: conclusion-2
% P:conclusion-06
\textit{[Insert the principal transfer finding, capture effect, and mobility-sensitivity conclusion.]}
% CHECKLIST-END: conclusion-2

% P:conclusion-07
The numerical study compares conservative simultaneous, shooting, finite-difference, and collocation formulations at common physical conditions and solution accuracy.
% CHECKLIST-BEGIN: conclusion-3
% P:conclusion-08
\textit{[Insert the method comparison, runtime effect, and refinement-supported accuracy.]}
% CHECKLIST-END: conclusion-3

% P:conclusion-synthesis
Together, the three studies relate absorber performance to the selected thermodynamic model, the local film representation, and the numerical work needed to resolve the coupled balances.
The operating comparison extends that interpretation to the specified neighborhood around Case 3C.
% CHECKLIST-BEGIN: conclusion-4
% P:conclusion-operating-finding
\textit{[Insert the supported operating tradeoff and the combined capture/temperature implication for model selection.]}
% CHECKLIST-END: conclusion-4


\subsection{Applicability and future work}\label{subsec:applicability}
% P:conclusion-09
The formulation applies to a steady, one-bed MEA absorber within the composition, temperature, and pressure domains of its thermodynamic and transport inputs.
Parameter uncertainty and extrapolation beyond the source ranges can affect speciation, transfer, and heat release together.
The equilibrium-manifold film assumes rapid local reaction and an isothermal diffusion path; conditions in which reaction or thermal relaxation competes with diffusion require finite-rate species and energy balances in the film.
Industrial gas mixtures containing additional reactive or volatile components require their species, reactions, properties, and interphase transfer terms.

% P:conclusion-10
The column balances and numerical formulations can be reused for DEA, MDEA, AMP, PZ, or blended solvents after defining the solvent-specific chemistry and constitutive inputs.
The extension begins with a reaction network, conserved components, and consistent activity and enthalpy reference states.
Pure-component, binary, association, ionic, and dielectric parameters are obtained from applicable sources or fitted to independent thermodynamic measurements, including VLE, speciation, and caloric data where available.
Transport properties and kinetic parameters are then specified for the solvent and packing.
Thermodynamic measurements excluded from fitting and independent absorber measurements evaluate the parameterization and the coupled column model, respectively.
A parameter-file substitution alone does not establish transferability to another solvent.

% P:conclusion-11
The specified input ranges of the Latin hypercube design define the domain of the operating comparison; they do not establish a global operating optimum.
Formal optimization additionally requires a stated objective and hydraulic, thermal, and process constraints.
Coupling the absorber to solvent regeneration would permit energy or economic objectives beyond capture alone.
These extensions should preserve the separation between property fitting, numerical verification, and experimental evaluation.
